% ============================================================
% chapters/00-abstract.tex -- תקציר מנהלים ומילות מפתח
% Executive summary: 250-350 words, one continuous block (booklet, sec. 2.2).
% Final version, written after all eight chapters (20.11.2026).
% Content only. The template decides where and how it is printed.
% Template must provide: environment `executivesummary`, command \keywords{...}
% ============================================================

\begin{executivesummary}
באוגוסט 410 בזזו הוויזיגותים את רומא, ומאז חוזרת השאלה אם קבלת הנצרות החלישה את הקיסרות הרומית. גיבון השיב בחיוב, ואילו רוב החוקרים בני זמננו מדגישים לחץ חיצוני, משבר כלכלי, אקלים ומגפות. לפי הטיעון המוצע בדוח, הוויכוח לא הוכרע בעיקר משום שמחקרים רבים בוחנים מקרה אחד, נשענים על דוגמאות נבחרות, ואינם מגדירים את המושגים שלהם באופן מדיד. דוח מכין זה מציע מחקר שבוחן את השאלה בהשוואה בין שלוש מדינות נוצריות שהגיעו לתוצאות שונות לגמרי: רומא המערבית, שנפלה ב־476; רומא המזרחית, שאימצה את אותה דת ושרדה עוד כאלף שנה; והממלכה הפרנקית, שהייתה נוצרית כמעט מהקמתה והתחזקה. שאלת המחקר היא אם הנצרות פעלה ככוח שמחליש או שמחזק מדינות בין המאה הרביעית למאה התשיעית, ומה קבע לאיזה כיוון היא פעלה. ההשערה המרכזית היא שהדת עצמה לא הכריעה, אלא מידת ההשתלבות בין הכנסייה לבין מוסדות המדינה. מבין שלוש חלופות שיטה נבחרה, במטריצת החלטה משוקללת, שיטה משולבת: מאגר נתונים השוואתי במקטעים של 25 שנה, שבו מדדי ההשתלבות מופקים מאלפי טקסטים עתיקים בעזרת מודלי שפה וועדת סוכני \en{AI}, ונבדקים מול תיוג של מומחה לתקופה על מדגם מבחן חתום, נפרד מקטעי הפיתוח. ניתוח רגישות הראה שהבחירה יציבה גם כשהמשקלים משתנים. המחקר נשען על כלים פתוחים קיימים – \en{Perseus}, \en{dMGH} ו־\en{Seshat} – ומוסיף להם מדידה שיטתית של היחסים בין הכנסייה למדינה. מפרט הדרישות קובע, בין היתר, שסיווגי המודל יתאימו לפחות ב־85\% לסיווג של מומחה, ושכל תוצאה תהיה ניתנת לשחזור. המחקר מתוכנן ל־12 חודשים בחמש חבילות עבודה, בתקציב תכנון של כ־111 אלף ש״ח, עם שלוש אבני דרך: תיקוף המודל בחודש החמישי, מאגר נתונים מלא בחודש השמיני, ופרסום פתוח של הקוד, הנתונים והממצאים בחודש ה־12. ארבעה סיכונים קריטיים זוהו, ובראשם הטיית המקורות שנכתבו ברובם בידי אנשי כנסייה וטעויות של מודלי השפה; לכל אחד מהם נקבעה תוכנית הפחתה, שלפי אומדן התכנון תוריד אותו מתחת לסף הקריטי. תוצרי המחקר יאפשרו לבחון את הוויכוח בן 1,600 השנים גם כשאלה אמפירית.
\end{executivesummary}

\keywords{נפילת רומא; התנצרות הקיסרות הרומית; הממלכה הפרנקית; יחסי דת ומדינה; יציבות מדינית; היסטוריה ממוחשבת; ניתוח טקסט בעזרת \en{AI}}
